\chapter{Convexité, tangentes et inégalités}
\section*{Exercice 15.1 -- Convexe mais décroissante}
La dérivée est $-e^{-x}<0$, donc la fonction décroît. La dérivée seconde est $e^{-x}>0$, donc elle est convexe sur $\R$. Ces deux propriétés portent sur des objets différents : valeur et pente.

\section*{Exercice 15.2 -- Inflexion}
$f''(x)=12x^2-4$ s'annule en $x=\pm1/\sqrt3$, change de signe à chacun et vaut négatif entre eux. Les deux points d'inflexion ont ordonnées $1/9-2/3=-5/9$.

\section*{Exercice 15.3 -- Tangente logarithmique}
Le logarithme est strictement concave ; sa courbe est sous la tangente en $1$, $y=x-1$. Pour l'égalité, $g(x)=x-1-\ln x$ a dérivée $(x-1)/x$ : son unique minimum est $g(1)=0$.

\section*{Exercice 15.4 -- Milieu}
Appliquer la définition de convexité de l'exponentielle à $t=1/2$. La stricte convexité donne égalité seulement si $a=b$. On retrouve aussi l'inégalité arithmético-géométrique pour les nombres $e^a,e^b$.

\section*{Exercice 15.5 -- Un faux critère}
Pour $f(x)=x^4$, $f''(x)=12x^2\ge0$ et $f''(0)=0$. La convexité ne change pas de part et d'autre de zéro ; il n'y a pas d'inflexion.

\section*{Exercice 15.6 -- Lorenz}
La somme vaut $10$. Les points sont $(0,0),(1/4,1/10),(1/2,1/5),(3/4,2/5),(1,1)$. Les pentes successives valent $0{,}4,0{,}4,0{,}8,2{,}4$, croissantes, donc la ligne polygonale est convexe.
